\clearpage
\VerdierPhysicalPageAnchor{0104}%
\begingroup
\small
\setlength{\abovedisplayskip}{3pt plus 2pt minus 2pt}%
\setlength{\belowdisplayskip}{3pt plus 2pt minus 2pt}%
\setlength{\jot}{2pt}%

\VerdierSemanticAnchor{ch1-3-4-3-zero-differential-complex}
Pour tout complexe $X$, désignons par $\overline{X}$ le complexe obtenu en
gardant les composants de $X$ et en annulant les différentielles. Il résulte
immédiatement de la définition de $H$ qu'on a l'égalité
$\overline{H(X)}=H(\overline{X})$. De même, pour tout morphisme de complexes
$f\colon X\longrightarrow Y$, désignons par
$\overline{f}\colon\overline{X}\longrightarrow\overline{Y}$ le morphisme obtenu
en conservant les composants de $f$. Il résulte alors de la définition du
foncteur cône que $c(\overline{f})=\overline{Y}\oplus\overline{X[1]}$, que le
morphisme $p(\overline{f})\colon\overline{Y}\longrightarrow c(\overline{f})$
est la première injection et que le morphisme
$q(\overline{f})\colon c(\overline{f})\longrightarrow\overline{X[1]}$ est la
deuxième projection.

\VerdierSemanticAnchor{ch1-3-4-3-cone-graded-splitting}
Le foncteur $H$ et le foncteur «\,$\overline{\phantom{X}}$\,» étant additifs,
on en déduit que
$\overline{H(c(f))}=\overline{H(Y)}\oplus\overline{H(X[1])}$, que
$\overline{H(p(f))}\colon\overline{H(Y)}\longrightarrow\overline{H(c(f))}$ est
la première injection et que
$\overline{H(q(f))}\colon\overline{H(c(f))}\longrightarrow\overline{H(X[1])}$
est la deuxième projection. Il en résulte que
\[
  H(c(f))^l=H(Y)^l\oplus H(X[1])^l,
\]

\VerdierSemanticAnchor{ch1-3-4-3-2}
\begin{equation}
  d_{H(c(f))}^{,l}=
  \begin{pmatrix}
    d_{H(Y)}^{,l} & \Phi^l\\
    0 & d_{H(X[1])}^{,l}
  \end{pmatrix},
  \tag{3.4.3.2}
\end{equation}

\VerdierSemanticAnchor{ch1-3-4-3-phi-map}
\[
  \Phi^l\colon H(X[1])^l\longrightarrow H(Y)^{l+1}.
\]
La différentielle de $H(c(f))$ s'obtient sous cette forme en exprimant
simplement que $H(p(f))$ et $H(q(f))$ sont des morphismes de complexes. Pour
calculer $\Phi^l$, il faut revenir à la définition du foncteur $H$ et, par
suite, à la définition du foncteur complexe simple associé (2.2.2.5). On a :

\VerdierSemanticAnchor{ch1-3-4-3-x-shift-components}
\[
  H(X[1])^l=
  \prod_{\substack{\sigma=\sum_i l_i e_i\\ \sum_i l_i=l}}
  F(X_1^{l_1},\ldots,X^{l_j+1},\ldots,X_n^{l_n}),
\]

\VerdierSemanticAnchor{ch1-3-4-3-3}
\begin{equation}
  H(Y)^{l+1}=
  \prod_{\substack{\sigma=\sum_i l_i e_i\\ \sum_i l_i=l}}
  F(X_1^{l_1},\ldots,Y^{l_j+1},\ldots,X_n^{l_n}),
  \tag{3.4.3.3}
\end{equation}

\VerdierSemanticAnchor{ch1-3-4-3-phi-components}
\[
  \Phi^l=
  \prod_{\substack{\sigma=\sum_i l_i e_i\\ \sum_i l_i=l}}
  \varepsilon(\sigma,j)
  F(X_1^{l_1},\ldots,f^{l_j+1},\ldots,X_n^{l_n}).
\]
\endgroup
